\section{Restrictions on finite-cover eightfold constructions}
\label{sec:finite-cover-eightfolds}
This section bounds the period images available from low-rank finite coefficient systems on curves. It restricts a method for constructing detectors; it gives no lower bound for the degrees of arbitrary cycles.

The reflection classification rests on Cohen's published classification and Malle's character-field tables. A recent unrefereed treatment by lelong\_ridge and cm\_symmetry\_scout \cite[Theorems~A--B and \S\S2,5]{ReflectionSources} classifies the primitive reflection possibilities for eightfold images of dimension at least fifteen and proves a primitive reflection-free coefficient-rank bound of twenty-three for that threshold. The two statements below use the sixteen-dimensional threshold for dominance. We give the induction and dimension arguments explicitly, including the case in which a complex block has no descent to the original character field.

\begin{theorem}[Gaussian reflection obstruction]\label{thm:gaussian-reflections}
Let $K=\Q(i)$. Consider a family of finite covers of smooth complex
curves and an absolutely irreducible finite coefficient representation
$\rho:G\to\operatorname{GL}_d(K)$ with exact character field
$\Q(\operatorname{tr}\rho)=K$. Suppose that its multiplicity in first
cohomology has dimension eight over $K$, and that its faithful coefficient
image contains a complex reflection. Then the corresponding period image
has complex dimension at most $15$. In particular it does not dominate
the signature-$(4,4)$ unitary domain of dimension $16$.
\end{theorem}

\begin{proof}
We give the reductions, including the normalizer and field-of-definition
steps. The classification inputs are the published results of Cohen
\cite{CohenReflection,CohenCorrection} and the reflection character fields
in Malle's tables \cite{MalleFields}. No unrefereed algebraicity theorem
is used.

\paragraph{The dimension bound and polarized induction.}
Let $V$ be a nontrivial irreducible finite coefficient system of rank $r$
on a genus-$\gamma$ curve, with genuine local monodromies $g_1,\ldots,g_b$.
Put $c_j=\operatorname{rank}(g_j-1)$. The Euler characteristic of its middle extension,
together with the absence of invariants in $V$ and its dual, gives
\[
 \dim_K H^1(Y,j_*V)=2r(\gamma-1)+\sum_{j=1}^b c_j.
 \label{eq:curve-euler}
\]
For rank-eight multiplicity and $r\ge2$, the corresponding Hurwitz
dimension is consequently at most
\[
 3\gamma-3+b\le8+(3-2r)(\gamma-1).
 \label{eq:curve-dimension}
\]
Positive-genus bases give a bound of eight; a genus-zero base gives
$2r+5$. For a nontrivial rank-one system, instead,
\[
 8=2\gamma-2+b,
 \qquad 3\gamma-3+b=\gamma+7\le12.
 \label{eq:rank-one-curve-dimension}
\]
This includes the unramified case. Fixed finite cyclic monodromy adds
only finitely many choices of root line bundles over a marked curve and
branch divisor, so the rank-one estimate bounds its period image. The middle-extension and
variation conventions agree with \cite[Proposition~1.1, Lemma~1.2,
Remark~1.3 and Theorem~2.5]{DettweilerWewers}.

If $\rho=\operatorname{Ind}_H^G\sigma$ over $K$, Shapiro identifies its multiplicity
cohomology with that of $\sigma$ on the intermediate curve
$C/H$. With $H_0=\ker\sigma$, the required faithful cover is
$C/H_0\to C/H$. This identification extends across branch points:
the stalks are local inertia invariants on both sides. It respects the
Hodge structure and the polarization up to the fixed nonzero sheet-trace
factor, since it is induced by the finite covering correspondences.
Thus the period map factors through this intermediate-cover family.
Forgetting additional cover data or points with trivial coefficient
monodromy does not enlarge its period image.

\paragraph{Reflection blocks are defined over $K$.}
Choose a positive invariant Hermitian form with entries in $K$ by
averaging. The root line $\operatorname{im}(s-1)$ of every reflection
$s$ is defined over $K$. Join two root lines when they are not orthogonal.
The connected components span mutually orthogonal $K$-subspaces $V_a$.
Their sum is nonzero and $G$-stable, hence is all of $V$.
The reflection normal closure is the product of its irreducible
reflection groups on these subspaces. The full group permutes the
blocks transitively. If $H$ stabilizes one block and $\sigma$ is its
action, then
\[
 \rho\simeq\operatorname{Ind}_H^G\sigma\quad\hbox{over }K,
 \qquad
 K\subset\Q(\operatorname{tr}\sigma)\subset K.
\]
The first inclusion follows from the induced-character formula. Thus
the faithful block stabilizer also has exact character field $K$.
It suffices, by polarized Shapiro, to treat this stabilizer and its
irreducible reflection subgroup $W$.

\paragraph{The classification and its low-rank cases.}
The groups that can occur for $W$ are rank-one groups, imprimitive
$G(m,p,r)$, rational real Weyl groups, and the nonreal primitive groups
$G_8,G_{29},G_{31}$, of ranks $2,4,4$ respectively. This is the
character-field restriction of Cohen's classification and Malle's
Tables~8.1--8.2; the omissions explained after those tables are
$G_4,G_{25}$ over the Eisenstein field and the real groups
$G_{28}=F_4$, $G_{35}=E_6$. Noncrystallographic real fields do not
embed in $K$.

For an imprimitive group of rank at least three, its diagonal elements
with entries $(\zeta_m,\zeta_m^{-1},1,\ldots)$ force
$\zeta_m+\zeta_m^{-1}\in\Q$. Hence $m\in\{1,2,3,4,6\}$.
The determinant-one diagonal element with entries
$(\zeta_m,\zeta_m,\zeta_m^{-2},1,\ldots)$ excludes $3$ and $6$.
The irreducible imprimitive cases therefore have $m=2$ or $4$.
All blocks of rank at most four have period dimension at most $13$
by the dimension estimate, with no restriction on their possible normalizer extensions.
This covers in particular triality in rank four.

\paragraph{High-rank monomial blocks.}
Cohen's Lemma~(2.7) gives a unique system of coordinate lines for
$G(m,p,r)$ when $r\ge5$. Every linear normalizer therefore preserves
that system. The reflection representation is conjugate over $K$ to
its usual monomial model over $K$: the intertwining equations have a
one-dimensional solution space over $K$, and a nonzero intertwiner is
invertible by absolute irreducibility. Thus its individual coordinate
lines, not merely their unordered collection, are defined over $K$.
The full block stabilizer is monomial on these lines.

For a line stabilizer $L$, its action is a finite character
$\psi:L\to K^\times$, and $\sigma=\operatorname{Ind}_L^H\psi$.
Exactness of the character field in the induced-character identity forces
$\Q(\psi)=K$, so its image is $\mu_4$. Polarized Shapiro realizes
the entire rank-eight multiplicity as a rank-one system on an
intermediate curve. The rank-one estimate bounds the period image by $12$.
Rank-one reflection blocks are covered by the same argument.

\paragraph{The rational high-rank normalizers.}
The remaining possibilities are $W(A_r)$ and $W(E_6),W(E_7),W(E_8)$.
In their rational reflection models, a finite-order $K$-matrix $g$
normalizing $W$ has the form $g=cA$ with $c\in K^\times$ and
$A\in\operatorname{GL}_r(\Q)$: solve the rational intertwining equations for the
automorphism of $W$ induced by $g$ and use absolute irreducibility.
The matrix $A$ is a rational similarity for the invariant positive
form and permutes root lines. All roots have the same length. Sending
one rational root to a rational multiple of another shows that its
similarity factor is a rational square. After dividing by that
rational multiple, $A$ is an orthogonal root-system automorphism.
It has finite order, so the remaining scalar is in $\mu_4$.

A root-system automorphism is a Weyl element followed by a Dynkin
diagram automorphism, by moving its chamber back to a fixed chamber.
For $A_r$ and $E_6$ the extra diagram involution is $-w_0$;
$E_7,E_8$ have no diagram automorphism. Hence the finite $K$-normalizer
is contained in $\mu_4W$. A subgroup containing $W$ and having exact
character field $K$ contains a nonreal coset, hence contains $iI$.
The faithful block stabilizer is therefore exactly $\mu_4W$.

For $E_6,E_7,E_8$ we may assume genus zero. In a generating
product-one branch tuple, at least two entries are outside the real
subgroup. Each is $i$ times a real orthogonal matrix, so its fixed
space has dimension at most $\lfloor r/2\rfloor$: the relevant
nonreal eigenvalues of the real matrix occur in conjugate pairs.
Each such entry moves at least $\lceil r/2\rceil$ dimensions.
Using the Euler formula gives
\[
 b\le2r+8-2(\lceil r/2\rceil-1).
\]
After subtracting three for automorphisms of the sphere, the period
dimensions are at most $13,13,15$ respectively.

\paragraph{All ranks in type $A$.}
It remains to avoid a rank-dependent gap in type $A$.
Assume $r\ge5$ and put $n=r+1$. Here $-I\notin W(A_r)$, so the
normalizer is the direct product $\mu_4\times S_n$, acting by
$\rho=\chi\otimes\operatorname{Std}_n$, with $\chi$ the scalar character.
We may again assume a genus-zero base. Let $b_\chi$ count branch
entries with nontrivial scalar phase. A phase $\pm i$ has moving
rank at least
\[
 a=r-\lfloor n/4\rfloor,
\]
because each permutation cycle divisible by four supplies at most
one eigenvector of the required eigenvalue. A phase $-1$ has moving
rank at least $c=r-\lfloor n/2\rfloor$. There are at least two
odd-phase entries, since the phases generate $\mu_4$ and their
product is one.

If $b_\chi\ge3$, choose two odd-phase entries and a third nontrivial
phase. All remaining genuine entries move at least one dimension.
Consequently
\[
 2r+8\ge2a+c+(b-3),\qquad
 b\le11-r+2\lfloor n/4\rfloor+\lfloor n/2\rfloor\le12.
\]
The period dimension is at most nine.

If $b_\chi=2$, the two phases are $i,-i$, and
$H^1(\mathbb P^1,j_*\chi)=0$ by the Euler formula. The representation identity
\[
 \chi\otimes\operatorname{Perm}_n
  =(\chi\otimes\operatorname{Std}_n)\oplus\chi
  =\operatorname{Ind}_{\mu_4\times S_{n-1}}^{\mu_4\times S_n}\psi
\]
therefore identifies the entire rank-eight cohomology in question
with the rank-one system $\psi$ on the degree-$n$
intermediate curve. There is no additional cohomology summand,
because the discarded $\chi$-summand vanishes. Its period dimension
is at most twelve by the rank-one estimate. These alternatives exhaust type $A$.

All reflection blocks and their finite compatible normalizers have
now been covered. Their bounds are at most fifteen, proving the theorem.
\end{proof}

\begin{theorem}[A coefficient-rank threshold for dominance]\label{thm:finite-cover-rank}
Let $V$ be a faithful irreducible complex representation of a finite group $G$, with nonrational character. Consider a Hurwitz family of $G$-covers whose $V$-multiplicity in first cohomology has complex dimension eight and Hodge numbers $(4,4)$. Suppose that every genuine branch monodromy has moving rank at least two. If its period image has dimension sixteen, then
\[
 \dim_{\C}V\ge26.
\]
The hypothesis concerns the branch monodromies; the group may contain reflections absent from their tuple. No descent of a complex imprimitivity system to a number field is assumed.
\end{theorem}

\begin{proof}
We first justify induction over $\C$, since it is needed even when the original coefficient representation has a smaller field of definition. If $V=\operatorname{Ind}_H^G W$ and $C$ is the covering curve, Frobenius reciprocity gives
\[
 \operatorname{Hom}_G(V,H^1(C,\C))
 \simeq\operatorname{Hom}_H(W,H^1(C,\C)).
\]
This is an isomorphism of variations of complex Hodge structures. To see the extension across the branch points, put $X=C/H$ and take the $W$-multiplicity in the finite pushforward of the constant sheaf from $C$ to $X$. A stalk at a branch point is the corresponding inertia-invariant space. Finite proper pushforward has no higher direct images and taking a finite-group isotypic part is exact, so this is the middle extension of the coefficient local system on the punctured curve. The maps on cohomology commute with the Hodge decomposition, since they come from finite sums of deck transformations.

The inducing representation $W$ is irreducible: otherwise induction of a proper subrepresentation would contradict irreducibility of $V$. It is nontrivial, since induction of the trivial representation contains constants. Its variation depends only on the finite cover $C/\ker(W)\to X$. If $r=\dim W$, $\gamma=g(X)$ and $b$ is its number of genuine coefficient branch points, the Euler formula is
\begin{equation}\label{eq:complex-cover-budget}
 8=2r(\gamma-1)+\sum_{j=1}^{b}c_j,
 \qquad c_j=r-\dim W^{g_j}\ge1.
\end{equation}
For fixed finite monodromy data the covers have finite fibres over the moduli of the marked base curve: its fundamental group has finitely many homomorphisms to the fixed finite group. Consequently the period-image dimension is at most $3\gamma-3+b$. For $r=1$ this is $\gamma+7\le12$. For $r\ge2$ it is at most
\begin{equation}\label{eq:complex-shapiro-bound}
 2r+5+(3-2r)\gamma\le2r+5.
\end{equation}
There is no unstable unmarked genus-one exception, because such a nontrivial finite system cannot have eight-dimensional cohomology. On a simply connected parameter neighborhood the Shapiro isomorphism is constant on the marked local systems. The filtration at one complex embedding determines the unitary period point, with its conjugate filtration determined by it. Thus forgetting the additional cover data does bound the original period image by \eqref{eq:complex-shapiro-bound}, without asserting number-field descent of $W$ or of a polarization on an intermediate abelian variety.

Assume for a contradiction that $d=\dim V\le24$ and the image has dimension sixteen. Rank one has already been excluded. For $d\ge2$ a positive-genus base gives image dimension at most eight, so the base is the sphere. Erase identity branch monodromies. The resulting product-one generating tuple has $b\ge19$, and
\begin{equation}\label{eq:rank25-budget}
 2d+8=\sum_{j=1}^{b}c_j\ge2b\ge38.
\end{equation}
Thus $15\le d\le24$. Some entry has moving rank two, because otherwise the sum would be at least $3b\ge57$, whereas $2d+8\le56$.

Suppose first that $G$ is imprimitive, with $t\ge2$ blocks of common dimension $r$ and $d=rt$. A block stabilizer gives an inducing representation of dimension $r$. Equation~\eqref{eq:complex-shapiro-bound} excludes $r\le5$, so $r\ge6$. Let $\sigma_j\in S_t$ be the permutation of blocks induced by the $j$th branch entry, and put $I_j=t-\#\{\text{cycles of }\sigma_j\}$. An operator cycling $\ell$ blocks has fixed space of dimension at most $r$, since a fixed vector is determined by its entry in one block. Hence $c_j\ge rI_j$. The block monodromy is transitive, and Riemann--Hurwitz for its degree-$t$ cover of the sphere gives
\[
 \sum_j I_j=2(t-1)+2\gamma'\ge2(t-1).
\]
If $J=\{j:I_j>0\}$, then $|J|\le\sum I_j$, while the other branch entries have rank at least two. Therefore
\[
 \sum c_j\ge2(b-|J|)+r\sum I_j
       \ge38+2(r-2)(t-1).
\]
Together with $d=rt$ and \eqref{eq:rank25-budget}, this implies $r+2t\ge17$. But $rt\le24$, $r\ge6$ and $t\ge2$ permit only $t=2,3,4$, for which the respective maxima of $r+2t$ are $16,14,14$. This contradiction excludes imprimitivity.

Next suppose that a tensor decomposition $V=A\otimes B$ is preserved, with $2\le a=\dim A\le s=\dim B$. For a finite-order operator $x\otimes y$, the two factors are diagonalizable up to a scalar: suitable powers of each are scalar. If $x$ is nonscalar, its eigenspaces have dimensions at most $a-1$, so the fixed space of $x\otimes y$ has dimension at most $(a-1)s$. Its moving rank is at least $s$. Since $d\ge15$, we have $s>2$, so a moving-rank-two element must have scalar $x$. Its remaining moving rank is a positive multiple of $a$, forcing $a=2$. Thus $d=2s$, with $s\le12$.

The projective action of $G$ on $A$ is nontrivial, since otherwise every subspace $\ell\otimes B$ would be invariant. A product-one generating tuple must therefore contain at least two entries nonscalar on $A$. They each have moving rank at least $s$; the other entries have rank at least two. It follows that
\[
 4s+8=\sum c_j\ge2s+2(b-2)\ge2s+34,
\]
or $s\ge13$, a contradiction. Here tensor indecomposability means that no such decomposition with both factors preserved exists; permutations of equal tensor factors are included in the convention of the classification theorem used next.

The published Theorem~7.1 and Remark~(3) of Guralnick--Saxl \cite[pp.~548--549]{GS} now apply: a finite irreducible primitive tensor-indecomposable complex group of degree greater than ten containing a moving-rank-two element is, up to scalars, the deleted permutation representation of the alternating or symmetric group. The other natural large-degree cases in that theorem have defining positive characteristic, and its exceptional characteristic-zero degrees are at most ten. Thus every element of $G$ has the form
\[
 g=\lambda\operatorname{Std}(\pi),\qquad\pi\in S_{d+1}.
\]
The scalar is unique, since no nonidentity permutation acts as a scalar on the deleted representation in these degrees. It therefore defines a character of $G$. If this character were trivial, the representation would have rational character. There are consequently at least two branch entries with $\lambda\ne1$.

A nonidentity eigenvalue of a permutation on $d+1$ letters has multiplicity at most $\lfloor(d+1)/2\rfloor$: each contributing cycle has length at least two and contributes one eigenvector. Removing the constant line does not change such an eigenspace. Each of the two nontrivial scalar entries thus has moving rank at least $\lfloor d/2\rfloor$. The branch budget gives
\[
 2d+8\ge2\lfloor d/2\rfloor+2(b-2)
       \ge2\lfloor d/2\rfloor+34.
\]
For even $d$ this forces $d\ge26$, and for odd $d$ it forces $d\ge25$. Both contradict $d\le24$. It remains to exclude $d=25$.

\paragraph{The degree-twenty-five case.}
The base is again the sphere, $b\ge19$, and $\sum c_j=58$. If the
representation is imprimitive, its block rank is either one or five.
The Shapiro bounds above give period-image dimension at most twelve
or fifteen, respectively. Thus $G$ is primitive. We show that a
primitive finite group of degree twenty-five with an element moving
at most four dimensions lies in
$\C^*\operatorname{Std}(S_{26})$.

First consider a preserved tensor structure. Its only nontrivial
factorization is $5\cdot5$. A nonscalar pure tensor has every
 eigenspace of dimension at most $4\cdot5=20$, by the eigenspace
argument above. A factor-interchanging operator has the form
$(A\otimes B)\tau$, with $\tau$ interchanging the factors. It is
conjugate to $(I\otimes C)\tau$, $C=BA$, whose square is
$C\otimes C$. Finite order makes $C$ semisimple. In a diagonal basis
there are five diagonal lines and ten two-dimensional off-diagonal
blocks. Each block contributes at most one vector to any fixed
 eigenspace. Its eigenspaces therefore have dimension at most fifteen.
Neither tensor case contains an element moving at most four dimensions.

The structural reduction in the proof of
\cite[Theorem~7.1, p.~549]{GS} applies to a primitive group with no
such tensor structure. If its Fitting subgroup is noncentral, it is
of extraspecial type, and four conjugates of a nonscalar element
 generate a subgroup containing a nontrivial scalar. Their fixed
spaces have zero intersection, so that element moves at least
$25/4$ dimensions. Otherwise there is a quasisimple normal subgroup
$E$ acting irreducibly on $V$. Indeed a nontrivial multiplicity space
or more than one commuting simple factor would give a tensor
structure. This representation of the matrix group $E$ is faithful.

Hiss--Malle's characteristic-zero classification
\cite[Theorem~1 and Table~2, pp.~22 and 31]{HissMalle}, with its
complete corrected table \cite[pp.~95--96 and 102]{HissMalleCorrection},
gives exactly the following degree-twenty-five possibilities:
\[
 E=A_{26},\quad \operatorname{PSL}_2(25),\quad
 \operatorname{PSL}_2(49),\quad \operatorname{PSp}_4(7).
\]
The alternating representation is deleted permutation. The two
$\operatorname{PSL}_2$ degrees are $q$ and $(q+1)/2$; the only other
ordinary row is $\operatorname{PSp}_4(7)$. The degree-twenty-five row
for ${}^3D_4(2)$ is in characteristic three and does not apply.

Here is the generation estimate, including scalar and outer lifts.
For a nonscalar $g$ with a nonzero fixed space, let $m$ be its
projective order, choose a prime $p\mid m$, and put $x=g^{m/p}$.
Then $x$ has projective order $p$. The scalar $x^p$ fixes every fixed
vector of $g$, so it is the identity. Thus $x$ has prime order and
moves no more dimensions than $g$. If $\alpha$ conjugates of its
induced automorphism generate a group containing $E/Z(E)$, the
intersection of the corresponding fixed spaces on $V$ is zero.
Otherwise their projective group would preserve a line, contradicting
irreducibility of $E$. Consequently
\[
 25\le\alpha\operatorname{rank}(x-I).
\]
For both $\operatorname{PSL}_2$ groups,
\cite[Lemma~3.1, pp.~527--529]{GS} gives $\alpha\le4$, including
field and diagonal automorphisms. For $\operatorname{PSp}_4(7)$,
\cite[Theorem~4.3, p.~533]{GS} gives $\alpha\le5$; its transvection
exception has the smaller value four. These bounds exclude an
 element moving at most four dimensions. The automorphism is
nontrivial, since an element centralizing the irreducible $E$ is
scalar, and a nonidentity scalar has no fixed vectors.

It remains that $E=\operatorname{Std}(A_{26})$. Its automorphism group
is $S_{26}$, and Schur's lemma then gives the full normalizer
$\C^*\operatorname{Std}(S_{26})$. This proves the claimed normalizer
assertion. In the branch tuple some entry moves at most three
 dimensions, since otherwise $58\ge4b\ge76$; hence it applies here.

Write each branch entry as $\lambda_j\operatorname{Std}(\pi_j)$.
A scalar phase outside $\{1,-1\}$ has finite order $k\ge3$, and a
fixed vector requires a cycle of $\pi_j$ whose length is divisible
by $k$. Therefore its moving rank is at least
\[
 25-\lfloor26/k\rfloor\ge17.
\]
If every scalar phase belonged to $\{1,-1\}$, the generated group
would have rational character. At least one phase is outside this
set; the product-one relation forces at least two. The remaining
branches each have moving rank at least two, giving
\[
 58=\sum c_j\ge2\cdot17+2(b-2)\ge68,
\]
a contradiction. This excludes degree twenty-five and proves the
 stated threshold.
\end{proof}

\begin{corollary}\label{cor:gaussian-cover-rank}
A dominant eightfold finite-cover source with an absolutely irreducible coefficient representation defined over $\Q(i)$ and having exact character field $\Q(i)$ has coefficient rank at least twenty-six.
\end{corollary}
\begin{proof}
Theorem~\ref{thm:gaussian-reflections} excludes a reflection in the faithful coefficient image. Every genuine branch entry therefore has moving rank at least two, and Theorem~\ref{thm:finite-cover-rank} applies.
\end{proof}

The degree-twenty-five profile with eighteen moving ranks three and
one moving rank four satisfies the Euler budget, but the proof shows
that it cannot occur in a dominant nonrational source. The distinction
between containing and being generated by small-moving elements matters:
$\mu_4\operatorname{Std}(S_{26})$ has nonrational character and contains
rank-three involutions, whereas all its rank-three involutions are
ordinary triple transpositions and generate a rational subgroup. A
positive construction in a larger degree still needs a finite group
with the correct character field, a split polarization, Hodge numbers
$(4,4)$ and a period differential of rank sixteen.



\subsection{Scalar-standard coefficients in every degree}

The alternating-group exception can be excluded without a degree
bound: its permutation realization forces the variation to come
from a small pointed curve.

\begin{theorem}\label{thm:scalar-standard-all}
Let a finite-cover source over $\mathbf P^1$ have an irreducible
complex coefficient representation $\rho$ of degree $d\ge26$,
compact cohomology multiplicity eight, and image contained in
$\C^*\operatorname{Std}(S_{d+1})$. Its period image has dimension
at most fifteen. This assertion does not require a particular
character field or a lower bound on the branch moving ranks.
Together with Theorem~\ref{thm:finite-cover-rank}, it excludes
scalar-standard coefficients in every degree for a dominant
nonrational source whose branch moving ranks are at least two.
\end{theorem}

\begin{proof}
Put $N=d+1$. The scalar and permutation factors are unique, so
\[
 \rho=\kappa\otimes\operatorname{Std}(\pi),\qquad
 \pi:G\longrightarrow S_N,\quad\kappa:G\longrightarrow\C^*
\]
are homomorphisms. Irreducibility makes the permutation action
transitive. Let $m$ count the branch places with nontrivial scalar
phase. At any such place, a fixed vector of $\rho$ comes from a
nontrivial eigenvalue of the permutation. If the phase has order
$k\ge2$, at most $\lfloor N/k\rfloor$ permutation cycles contribute,
so its moving rank is at least
\[
 d-\lfloor N/2\rfloor=\lfloor d/2\rfloor.
\]
The compact Euler formula therefore gives
\[
 m\lfloor d/2\rfloor\le2d+8.
\]
Since $5\lfloor d/2\rfloor>2d+8$ for $d\ge26$, we obtain $m\le4$.

Let $q:Y\to\mathbf P^1$ be the connected degree-$N$ intermediate
cover defined by the permutation action. Let $\mathcal L$ be the
rank-one finite local system of $\kappa$, and let $\mathcal L_Y$
be its pullback. Compact Shapiro and
$\operatorname{Perm}_N=\mathbf1\oplus\operatorname{Std}$ give
\begin{equation}\label{eq:standard-shapiro}
 H^1(Y,j_*\mathcal L_Y)
 =H^1(\mathbf P^1,j_*\mathcal L)\oplus H^1_\rho.
\end{equation}
Here points with trivial upstairs inertia are filled in. To check
this convention at a branch point, a permutation cycle of length
$e$ contributes an invariant stalk exactly when the scalar phase
to the power $e$ is one; this is also exactly when the upstairs
rank-one system extends across that point. The finite-pushforward
projection formula therefore identifies the middle extensions as
well as their unramified restrictions. The equality is compatible
with Hodge structures and Gauss--Manin connections.

If $\kappa$ is nontrivial, its compact cohomology on the sphere has
dimension $m-2$, so the left side of
\eqref{eq:standard-shapiro} has dimension
\[
 n=m+6\le10.
\]
Suppose first that $\mathcal L_Y$ is nontrivial. If $Y$ has genus
$h$ and the rank-one system has $r$ genuinely ramified points, then
\[
 n=2h-2+r.
\]
The finite-order local system adds only discrete choices to an
$r$-pointed genus-$h$ curve. For $h\ge2$ its moduli dimension is
\[
 3h-3+r=n+h-1\le\tfrac32 n\le15.
\]
For $h=1$ it is $r=n\le10$, and for $h=0$ it is
$r-3=n-1\le9$. If $\mathcal L_Y$ is trivial, then $2h=n\le10$,
and the unpointed curve moduli dimension is at most twelve.
Finally, if $\kappa$ itself is trivial, the first summand in
\eqref{eq:standard-shapiro} vanishes and $Y$ has genus four,
giving dimension at most nine.

These bounds control the selected summand, not merely the full
left side of \eqref{eq:standard-shapiro}. Work locally on a simply
connected marked neighborhood of the cover family. That summand
is a fixed flat subspace, and its Hodge filtration is the intersection
with the ambient filtration. A tangent vector fixing the pointed
curve and its finite local system therefore fixes the summand's
filtration as well. Additional motions of the map to $\mathbf P^1$,
or of points whose upstairs inertia is trivial, cannot increase
its period rank. This proves the asserted bound.
\end{proof}

\subsection{Ordered tensor factors require degree at least forty-eight}

A fixed tensor factor imposes a further branch budget in every degree,
independently of the character field.

\begin{proposition}\label{prop:ordered-tensor48}
Suppose a finite irreducible coefficient representation
$V=A\otimes B$, with $2\le a=\dim A\le b=\dim B$, preserves this
ordered tensor decomposition.  If a sphere-cover family has compact
cohomology multiplicity eight and period image of dimension sixteen,
then
\[
 2ab+8\ge(a+1)b+18a-a^2,\qquad ab\ge48.
\]
The assertion does not cover elements permuting tensor factors.
\end{proposition}

\begin{proof}
Let $g_1,\ldots,g_N$ be the nonidentity branch matrices.  They generate
the coefficient image, have product one, and $N\ge19$.
Choose a representative $x_i$ of their projective action on $A$,
scaled so that a largest eigenspace has eigenvalue one.  Put
\[
 \ell_i=a-\dim\ker(x_i-1),\qquad L=\sum_i\ell_i.
\]
Finite projective order makes each $x_i$ diagonalizable.  The product
of these representatives is scalar.  The projective action on $A$
is irreducible, since a proper invariant subspace would tensor with
$B$ to a proper invariant subspace of $V$.

Delete the scalar $x_i$.  At least two entries remain: a single
nonscalar matrix cannot have scalar product.  If $L\le a$, the
intersection of the fixed spaces of all but the last remaining
entry has dimension at least the positive moving rank of that last
entry.  It is nonzero and proper.  All preceding entries fix it
pointwise; the scalar product relation makes the last entry scalar
on it.  This contradicts irreducibility.  Thus $L\ge a+1$.

Write $c_i=\operatorname{rank}(g_i-1)$.  If $\ell_i>0$, an eigenspace
of the $B$ component matches at most one eigenspace of $x_i$, of
dimension at most $a-\ell_i$.  Hence $c_i\ge b\ell_i$.
If $\ell_i=0$, the nonidentity matrix $g_i$ moves at least one whole
$A$ factor, so $c_i\ge a$.  With $r$ the number of positive $\ell_i$,
we have $r\le L$, and therefore
\[
 \sum_i c_i\ge bL+a(N-r)
    \ge aN+(b-a)L
    \ge19a+(b-a)(a+1).
\]
The Euler equality $\sum_i c_i=2ab+8$ proves the first inequality.
Rearranging it gives $(a-1)b\ge18a-a^2-8$.
For $a=2,3,4,5,6$, the least possible $b$ is respectively
$24,19,16,15,13$, giving products $48,57,64,75,78$.
For $a\ge7$, the product is already at least $49$.
This proves the degree bound.  Equality at $(a,b)=(2,24)$ is only
a numerical possibility, with no group or cover existence implied.
\end{proof}

\subsection{An imprimitive Gaussian representation of degree ninety-nine}

A larger coefficient rank does not by itself remove the branch-cost
obstruction. The following example isolates the extra cost imposed by
its permutation action on blocks.

\begin{proposition}\label{prop:degree99-cover}
Let $G$ be the elementary-abelian extension $2^{10}.M_{11}$ identified
in the Character Table Library by
\texttt{P48/G1/L1/V2/ext2}, and let $\rho$ be either of its conjugate
faithful degree-$99$ characters with character field $\Q(i)$.
A sphere-cover source for $\rho$ with compact cohomology multiplicity
eight has at most seven nontrivial coefficient branch points. Its
period image therefore has dimension at most four. The assertion
also holds after finite scalar enlargement of the coefficient group.
\end{proposition}

\begin{proof}
The group construction and its table are given in
\cite[\S1.2-9]{CTblLib}. Write $N$ for the elementary abelian normal
subgroup of order $1024$. The seven $G$-classes in $N$, in increasing
library order, give the following restriction data:
\[
\begin{array}{c|rrrrrrr}
\text{class size}&1&55&330&396&66&165&11\\
\chi&99&63&27&-9&-9&-45&-81.
\end{array}
\]
The two conjugate characters have the same restriction, since these
values are real. Thus
\[
 \langle\chi|_N,\chi|_N\rangle_N
 =\frac{99^2+55\cdot63^2+330\cdot27^2
       +462\cdot9^2+165\cdot45^2+11\cdot81^2}{1024}
 =891.
\]
Clifford theory expresses this restriction as $m$ times an orbit of
$n$ distinct linear characters of $N$. Since $mn=99$ and
$m^2n=891$, we obtain $m=9$ and $n=11$. The corresponding eigenspaces
are eleven blocks of dimension nine. Their permutation action
factors through $M_{11}$, is transitive, and is faithful because
$M_{11}$ is simple and the action is nontrivial.

The degree-eleven transitive permutation character of $M_{11}$ is
\[
 (11,3,2,3,1,0,1,1,0,0).
\]
For clarity, its identification does not require choosing a block
labeling: the least nontrivial irreducible character degree is ten,
so a degree-eleven transitive permutation character is one plus a
degree-ten irreducible; exactly one such sum has nonnegative integral
values. The class power maps give permutation indices
\[
 (0,4,6,6,8,8,8,8,10,10).
\]
Equivalently, the number of permutation cycles is the multiplicity
of the eigenvalue one, which gives these indices by subtraction
from eleven. In particular every nonzero index is at least four.

Let $i_j$ be the block-permutation index of a coefficient branch
matrix and let $c_j$ be its moving rank. A permutation cycle on
$\ell$ blocks contributes at most nine fixed vectors and hence at
least $9(\ell-1)$ moving vectors. Therefore $c_j\ge9i_j$, with any
scalar phase allowed. If $i_j=0$, the matrix lies in $N$ up to a
scalar and acts as a scalar on each block. A nonidentity matrix
then moves at least one whole block, so $c_j\ge9$. This inequality
also follows from the first one when $i_j>0$. Since a nonzero
$i_j$ is at least four, the two bounds imply
\[
 c_j\ge9+\frac{27}{4}i_j.
\]

The branch permutations generate a transitive degree-eleven group
and have product one. Riemann--Hurwitz for the associated degree-eleven
cover gives $\sum i_j\ge20$. If there are $b$ coefficient branches,
\[
 206=2\cdot99+8=\sum_jc_j
     \ge9b+\frac{27}{4}\sum_ji_j
     \ge9b+135.
\]
Consequently $b\le7$, and the branch-configuration dimension is
$b-3\le4$. All estimates were independent of scalar phases, proving
also the last assertion.
\end{proof}

The numerical restriction data can be verified from the published
permutation-group construction, without accepting an abstract table
as a group-realization certificate. Computing that group's order gives
$8110080$; its solvable radical is elementary abelian of order $1024$;
and its computed ordinary character table has $52$ classes and agrees
with the cited table. The \emph{Gaussian finite-cover checks} in the
\href{\workbenchurl}{workbench} reproduce these exact calculations.
A scan of the $2678$ named tables in version~1.3.11 selects $942$
faithful exact-Gaussian characters of degrees $25$ through $250$.
Only the two characters of this proposition pass the initial necessary
inequality $19\nu\le2d+8$, where $\nu$ is the least nonzero projective
moving rank. This describes a finite search and its coverage; it is
not a classification of finite Gaussian linear groups. In particular,
scalar twists of rational representations need not occur in that list.

\subsection{A degree-one-hundred-ten coefficient over an imaginary quadratic field}

The same exact test over the other imaginary quadratic character
fields leaves one additional pair in that finite catalogue.  Its
branch costs already contradict the compact cohomology dimension.

\begin{proposition}\label{prop:degree110-cover}
Let $G$ be the group $2^{10}.M_{11}$ identified in the Character Table
Library by \texttt{P48/G1/L1/V1/ext2}, and let $\rho$ have either
faithful degree-$110$ character with exact character field
$\Q(\sqrt{-2})$.  No sphere-cover source for $\rho$ with compact
cohomology multiplicity eight has a sixteen-dimensional period image.
The assertion permits arbitrary finite scalar enlargement.
\end{proposition}

\begin{proof}
The ordinary character table and its power maps give least nonzero
moving rank ten, even after arbitrary scalar phases.  The Euler
budget is $2\cdot110+8=228$.  If there are at least nineteen branch
points, each must have moving rank at most $228-18\cdot10=48$.
Here is the complete list of class and scalar pairs with moving
rank between two and forty-eight; the conjugate characters have the
same list:
\[
\begin{array}{c|r|r}
 \text{class}&\text{scalar}&\text{moving rank}\\\hline
 2a&1&20\\
 2b&1&40\\
 2e&-1&30\\
 2f&-1&10
\end{array}
\]
For each class of order $s$, the power-map traces determine the
multiplicity of every $s$th root of unity by finite Fourier inversion.
A scalar phase can give a fixed vector only when it is the inverse
of one of those eigenvalues; every other phase has moving rank
$110$.  Enumerating these inverse phases gives exactly the four
rows displayed.  Thus all branch costs are divisible by ten, whereas
$228$ is not.  This is impossible before any generation or
polarization condition is imposed.

There is also a quotient check.  The nontrivial degree-ten character
of the table with index two is trivial on the scalar centre and on
all four displayed classes.  A product-one tuple generating the
projective image would have, in that character, total moving rank
at least twenty, by the nonnegativity of its compact first cohomology.
The four classes would give total moving rank zero.  This gives
the same exclusion independently of the divisibility test.
\end{proof}

For context, scanning all $2678$ named ordinary tables in version
1.3.11 selects $1592$ faithful characters of degree $26$ through
$250$ with exact imaginary quadratic character field other than
$\Q(\sqrt{-3})$.  They lie in $371$ named tables and nineteen
imaginary quadratic fields.  Only the two conjugate pairs in
Propositions~\ref{prop:degree99-cover} and
\ref{prop:degree110-cover} satisfy the necessary minimum-cost
inequality, even with arbitrary scalar phases allowed.
The \emph{quadratic finite-cover checks} in the
\href{\workbenchurl}{workbench} give the exact field tests,
power-map computations, branch costs and rational certificates.
These are computations in a specified finite catalogue; they do
not classify all finite groups, assert a Schur index, or construct
an integral coefficient lattice.

\subsection{Rational coefficient characters and arbitrary scalar phases}

There is a separate restriction when the coefficient character before
scalar enlargement is rational.  It permits rank-one moving parts and
all finite scalar phases.

\begin{lemma}\label{lem:rational-scalar-budget}
Let $\rho$ be a finite irreducible representation of degree $d>1$
with rational character.  Suppose that a product-one tuple
$M_j=\lambda_j\rho(g_j)$ generates an irreducible group with nonreal
character, has at least nineteen genuine branches, and has parabolic
first cohomology of dimension eight.  Let $\nu$ be the least positive
moving rank among all scalar multiples of elements of $\rho(G)$,
and let $\eta$ be the least moving rank with nonreal scalar phase.
Then
\[
 2\eta+17\nu\le 2d+8,\qquad \eta\ge\lceil d/2\rceil.
\]
\end{lemma}
\begin{proof}
A scalar in the rational-character image has rational trace $d\zeta$,
so it is $1$ or $-1$.  The product-one condition consequently makes
$\prod_j\lambda_j$ real.  All real phases would give rational
character on every generated word.  At least one phase is therefore
nonreal, and the real product forces at least two such phases.
The Euler formula gives $\sum_j\operatorname{rank}(M_j-1)=2d+8$.
Two terms cost at least $\eta$, and at least seventeen others cost
at least $\nu$.  Finally the rational character gives equal
multiplicities to conjugate eigenvalues.  A nonreal eigenspace has
dimension at most $\lfloor d/2\rfloor$, which proves the last bound.
\end{proof}

\begin{proposition}[A bounded-degree symmetric-family exclusion]
\label{prop:symmetric-scalar-exclusion}
Under the hypotheses of Lemma~\ref{lem:rational-scalar-budget},
$\rho$ cannot be a nonstandard irreducible representation of any
$S_s$ of degree $26\le d\le500$, or a scalar enlargement thereof.
Here nonstandard excludes the trivial, sign, standard and
sign-twisted standard representations.  This conclusion is a finite
exact computation together with the degree cutoff proved below.
\end{proposition}
\begin{proof}
For a cycle class of order $h$, rationality makes the multiplicity
$m_q$ of an individual primitive $q$th root independent of its choice.
Fourier inversion, grouping powers by their order, gives
\[
 m_q=\frac1h\sum_{r\mid h}\chi(g^{h/r})c_r(h/q),\qquad
 c_r(e)=\mu\left(\frac r{(r,e)}\right)
       \frac{\varphi(r)}{\varphi(r/(r,e))},\quad q\mid h.
\]
The second expression is the sum of the primitive $r$th roots to
the $e$th power, obtained by inclusion--exclusion over divisors.
Thus the real-phase costs are $d-m_1,d-m_2$, and the nonreal
minimum is $d-\max_{q>2}m_q$.  A phase whose inverse is absent
from the spectrum costs $d$ and needs no additional enumeration.
Zero cost is omitted from the genuine branches.

The exact Murnaghan--Nakayama calculation in the
\emph{rational scalar checks} of the
\href{\workbenchurl}{workbench} enumerates the $117$ nonstandard
partitions up to transposition with $5\le s\le33$ and degree
between $26$ and $500$.  Transposition tensors with sign, which
preserves all phase-cost minima.  For every cycle partition the
calculation checks nonnegative integral multiplicities and
\[
 \sum_{q\mid h}\varphi(q)m_q=d,\qquad
 \sum_{q\mid h}\mu(q)m_q=\chi(g).
\]
Every selected character has
$2\eta+17\nu-(2d+8)>0$.  The smallest excess is $76$,
at the partition $(7,2)$ of $9$, of degree $27$, with
$\nu=6$ and $\eta=18$.  For comparison, $(31,2)$ of $33$
has degree $495$, $\nu=30$, $\eta=330$ and excess $172$.
The stored rows give the degrees, minima and witnessing cycle
partitions; a separate integer hook-length calculation verifies
the completeness of the partition selection.

It remains to justify the cutoff for all $s$.  A nonstandard
partition of $s$ has first row and first column at most $s-2$.
For $s\ge6$ it has a removable corner leaving a nonstandard
partition of $s-1$.  If both lengths are at most $s-3$, any
corner works.  If the first row has length $s-2$, the remaining
two boxes form $(2)$ or $(1,1)$; removing the last box of the
first row works.  The column case is its transpose, and the two
exceptional bounds cannot hold simultaneously for $s\ge6$.
The restriction branching rule, or equivalently the decomposition
of standard tableaux by the corner containing their largest entry,
makes the degree at least that of this predecessor.
The independent exact hook-length enumeration of all $12310$
partitions of $34$ gives smallest nonstandard degree $527$,
attained only by $(32,2)$ and its transpose.  Repeated removal
therefore excludes degrees at most $500$ for every $s\ge34$.
For $s<5$ all degrees are below $26$.  This proves the assertion.
\end{proof}

The analogous named-table calculation in CTblLib version~1.3.11
reads all $2678$ ordinary table names and selects $12479$ faithful
rational-character rows of degrees $26$ through $500$, occurring
in $1365$ named tables.  Every row violates
Lemma~\ref{lem:rational-scalar-budget}, with smallest strict
excess $28$, even with arbitrary scalar phases and moving rank
one allowed.  Exact cyclic Fourier inversion is used in every
case, with dimension, trace, integrality and conjugation checks.
This last assertion has precisely the scope of that finite
catalogue.  It does not exhaust unlisted finite groups, tensor
products or proper irreducible projective subgroups.
